\chapter*{Resumen Trilingüe}
\label{sec:abstract}
\vspace*{-10mm}

\textbf{Resumen}

\textit{"En la última década, la industria ha llevado a cabo una transición acelerada desde arquitecturas centralizadas hacia sistemas de control distribuidos (DCS). Este nuevo estandar, impulsado por la digitalización y la convergencia de las Tecnologías de la Operación (OT) y de la Información (IT), permite crear redes de sistemas ciberfísicos (CPS) escalables sentando las bases de la Industria 4.0. Este cambio es críticamente notable en el sector energético, donde el modelo tradicional de generación de energía, basado en combustibles fósiles en centrales monolíticas, ha dado paso a un sistema más diversificado, sostenible y distribuido (parques eólicos, plantas solares etc.). En este contexto, los sistemas de supervision, control y adquisicion de datos (SCADA) se han consolidado ofreciendo soluciones de conectividad que garantizan la interoperabilidad, seguridad en los CPS. El trabajo de fin de master consiste en una solución SCADA para una red de generación eólica, el cual estará ubicado en una sala de control central (MCR).El desarrollo se fundamenta en la normativa internacional IEC 61400-25 que regula las comunicaciones en parques eólicos. Este sistema SCADA posibilita la gestión de alarmas y eventos, la elaboración de gráficas de tendencia y la ejecución de tareas de mantenimiento predictivo de las instalaciones desde un punto central, mejorando la eficiencia y productividad de toda la red. La memoria recoge además un manual de usuario y un banco de pruebas.}\par

% En Español: 
\textbf{Palabras clave:} Industria 4.0; Sistemas CiberFísicos; Parque Eólico; Energía; Subestaciones eléctricas; Energía sostenible
; OPC UA; Modbus TCP ; mantenimiento; turbinas.

\vspace*{5mm}
\textbf{Abstract}

\textit{In the last decade, the traditional model of power generation, based on centralized fossil fuels, has given way to a more diversified, decentralized, and sustainable system. As the new paradigm becomes established, information management problems arise due to the dispersion of systems. To address the logistics and management issues caused by this dispersion, the master's thesis proposes the design of a supervisory control and data acquisition system (SCADA) supported by the international standard for communication in wind power plants IEC 61400-25, to unify signals and enable the centralized and efficient control that the new distributed generation standard requires. This SCADA system enables alarm and event management, the creation of trend graphs, and the execution of predictive maintenance tasks for installations from a central point, improving the efficiency and productivity of the entire network.}\par

\textbf{Keywords:} xxx
\vspace*{5mm}
\textbf{Laburpena}

Azken hamarkadan, erregai fosil zentralizatuetan oinarritutako energia-sorkutza bide eman dio eredu deszentralizatu batera. Paradigma berria ezartzen den heinean, informazioaren kudeaketan arazoak sortuz joan dira. Master amaierako lanak, energia berriztagarri ekoizpen sare baten kontrol eta kudeaketa gela (MCR) diseinatzen du. Horrela, estazio monolitikoetan tradizionalki zegoen kontrolaren antzeko sistema bat ezartzen da eta kontrol banatuaren arazoari irtenbidea ematen zaio. Gelak gainbegiratzeko kontrol eta datu eskuratze sistema (SCADA) bat eta WIN Server sistema eragileekin bi zerbitzari erreserbatu ditu, eta 20 azpiestazio elektriko eta energia-sorkuntza sistemetako bakoitzarekin komunikatuta dago (122 aerosorgailu eta eguzki-parke), guztira 744 MW energia elektriko kudeatuz. Horrela, MCRk alarmen eta gertaeren kudeaketa, joeren grafikoen prestaketa eta instalazioen mantentze prediktiboko lanak egitea ahalbidetzen du. Azkenik, martxan jartze birtual bat gauzatzen du funtzionamendu egokia bermatzeko.

\textbf{Hitz gakoak:} xxx


